\section{Related work}
%-----------------------------------------
In this section, we comparatively discuss BaroVox with state-of-the-art works in the following two categories. 


{\textbf{Side-channel speech eavesdropping:}} Side-channel eavesdropping has been considerably studied in academic literature~\cite{michalevsky2014gyrophone, zhang2015accelword, anand2018speechless, anand2021spearphone, han2017pitchln, 7878599, marquardt2011sp,9517289, 9095984, cpams}. Several works, such as \cite{michalevsky2014gyrophone, zhang2015accelword, anand2018speechless, anand2021spearphone, han2017pitchln, marquardt2011sp, ba2020learning, wang2022mmphone, gao2022device, gao2022kite}, have investigated the vulnerabilities of motion sensor eavesdropping, most notably in mobile devices.
Sami et al. \cite{sami2020spying} devised a novel acoustic side-channel threat using the lidar sensors found in consumer-grade robot vacuums. 
Roy et al. \cite{roy2016listening} looked into the practicality of using a mobile device's vibration motor as a sound sensor. Nassi et al. \cite{nassi2022lamphone} exploit the vibrations of a hanging bulb inside a room to retrieve sound from desk lamp light bulbs through an optical side-channel attack. 
\cite{kwong2019hard,long2023side,nassi2021glowworm,nassi2023little,nassi2022little,wang2022wavesdropper,basak2022mmspy} are other works that focus on acoustic signal eavesdropping. 
These studies raised public awareness of the feasibility of recovering sound by analyzing non-acoustic data. Our work takes a unique approach to side-channel attacks by demonstrating the \textit{\textbf{first-ever use of this technique on pressure sensors}}, adding a new dimension to this field of study. 

\textbf{Attacks on pressure sensors:}
Tu et al. \cite{tu2021transduction} reveal threats of Electromagnetic interference spoofing attacks on tire pressure sensors to over/under-inflate car tire.
Barua et al. designed malicious music to create resonance in pressure sensors to fool the pressure sensor used in the RPM systems of a negative pressure room (NPR) to turn NPR's negative pressure into positive one~\cite{barua2022wolf}.   Our research distinguishes itself by highlighting the sensitivity of pressure sensors across a range of frequencies beyond the resonant frequency. Moreover, while both papers engage DPS sensors, they diverge significantly in focus and methodology. 
Barua et al. exploit resonant frequencies to manipulate pressure readings and create hazardous conditions --- an integrity attack. In contrast, our work reveals the potential of DPS as a covert channel for leaking acoustic information 
 --- a confidentiality attack. Additionally, we demonstrate the potential to extract sound signals from minor fluctuations in pressure readings, a new avenue for exploiting pressure sensors.

% \textbf{Attacks on other sensors:}